\maketitle

\renewcommand{\abstractname}{Abstract}
% \selectlanguage{english}
\begin{abstract}


\noindent Evidence that mass transit raises local income does not settle whether the gains reach the residents who already lived in the areas it serves, and most of it comes from advanced economies. This paper addresses this gap by studying the effect of the Brasília subway on local income. Brasília, a planned capital whose core is legally protected from densification, is an extreme case of a city where residents cannot move closer to jobs. Using census-tract data for 2000 and 2010 and an instrumental-variable strategy based on a planned and discarded subway alignment, we find that tracts exposed to the subway experienced income growth approximately 10 percent higher than comparable unexposed tracts. We find no evidence that this gain reflects physical restructuring of exposed areas or the replacement of lower-income residents by wealthier ones. In addition, both the share of household heads with positive income and the average income of those who already had it rise more in exposed tracts, consistent with a labor-market channel. Overall, the results indicate that, where moving closer to jobs is not an option, a subway can raise the income of the residents who stay.

\noindent
\end{abstract}

\keywords{Mass transit; Income; Brasília.}\

\jel{O18, R11, R23, R42.}

% \doublespacing
% \thispagestyle{empty}
\clearpage

\onehalfspacing
